% Study note: the linear algebra behind the identifiability analysis.
% Build:  latexmk -pdf linalg_note.tex   (PDF is gitignored)
\documentclass[11pt]{article}
\usepackage[margin=1in]{geometry}
\usepackage{amsmath,amssymb,amsthm}
\usepackage{booktabs}
\usepackage{xcolor}
\usepackage[colorlinks=true,allcolors=blue]{hyperref}
\setlength{\parskip}{4pt}
\setlength{\parindent}{0pt}
\newcommand{\Mat}[1]{\mathbf{#1}}
\newcommand{\A}{\Mat{A}}
\newcommand{\thet}{\boldsymbol{\theta}}
\newcommand{\Rw}{R^2_{\widehat W}}
\DeclareMathOperator{\rank}{rank}
\DeclareMathOperator{\cond}{cond}
\DeclareMathOperator{\tr}{tr}
\newcommand{\eps}{\varepsilon}
\newcommand{\src}[1]{{\footnotesize\textcolor{gray}{[#1]}}}
\newtheorem{theorem}{Theorem}
\newtheorem{remark}{Remark}

\title{Least squares, the Gram matrix, and the eigenvalue spectrum\\
\large What the identifiability analysis actually computes, and where it can mislead}
\author{Study note for submission 28595}
\date{}

\begin{document}
\maketitle
\vspace{-2em}

\begin{abstract}
\noindent Part~I derives the linear algebra from scratch: least squares, the normal
equations, the Gram matrix, the SVD, the eigenvalue spectrum, conditioning,
truncation, and the two perturbation results that make ``noise breaks
degeneracies'' a theorem rather than an intuition. Every equation is attributed.
Part~II lists every place these objects appear in \texttt{neurips.tex} and
\texttt{reply\_all.tex}, explains each line, and flags four places where the
wording does not match what the code computes --- including the one the Area
Chair objected to. Part~III resolves the apparent conflict between the
eigenvalue analysis and the least-squares recovery numbers.
\end{abstract}

\part{The tools}

\section{The problem and its normal equations}

Everything below concerns one fixed postsynaptic neuron $i$. The Flyvis ODE
\src{Lappalainen et al.\ 2024}
\begin{equation}
\tau_i \dot v_i = -v_i + V^{\mathrm{rest}}_i
  + \textstyle\sum_{j\in\mathcal N_i} W_{ij}\,\mathrm{ReLU}(v_j) + I_i
\label{eq:ode}
\end{equation}
is \emph{linear in its parameters} $\thet_i = (\tau_i, V^{\mathrm{rest}}_i,
W_{ij_1},\dots)$ once the trajectory is observed. Getting from
\eqref{eq:ode} to \eqref{eq:ls} is nothing but moving terms across the
equals sign. Collect on the left every term that contains an
\emph{unknown} and on the right every term that is \emph{known} once $v$
has been recorded:
\begin{equation}
\tau_i \dot v_i \;-\; V^{\mathrm{rest}}_i
  \;-\; \textstyle\sum_{j\in\mathcal N_i} W_{ij}\,\mathrm{ReLU}(v_j)
\;=\; -v_i + I_i .
\label{eq:ode-rearranged}
\end{equation}
Every unknown now appears exactly once, multiplied by a recorded quantity,
so the left-hand side is a dot product of a known row vector with $\thet_i$:
\begin{equation}
\underbrace{\big(\dot v_i,\ -1,\ -\mathrm{ReLU}(v_{j_1}),\ \dots,\
  -\mathrm{ReLU}(v_{j_{d_i}})\big)}_{\text{known: one row of }\A_i}
\cdot
\underbrace{\big(\tau_i,\ V^{\mathrm{rest}}_i,\ W_{ij_1},\ \dots,\
  W_{ij_{d_i}}\big)^{\!\top}}_{\textstyle \thet_i}
\;=\; \underbrace{-v_i + I_i}_{\text{known: one entry of }\Mat b_i} .
\label{eq:ode-dot}
\end{equation}
Writing \eqref{eq:ode-dot} out at each of the $T$ sampled times and stacking
those rows gives
\begin{equation}
\A_i\,\thet_i = \Mat b_i, \qquad \A_i \in \mathbb R^{T\times n},\quad n = 2 + d_i,
\label{eq:ls}
\end{equation}
with $d_i$ the in-degree: $T$ rows because one equation per timestep,
$n=2+d_i$ columns because one unknown per parameter. This is Suppl.\
Eq.~(6)--(7) of the submission. The point worth internalising: no
linearisation has occurred. The nonlinearity $\mathrm{ReLU}(v_j)$ sits in the
\emph{data} $\A_i$, not in the unknowns --- rearranging
\eqref{eq:ode}$\to$\eqref{eq:ode-rearranged} did not approximate anything, it
only sorted the terms.

Since $T \gg n$, \eqref{eq:ls} is overdetermined and we solve it in the
least-squares sense,
\begin{equation}
\hat\thet_i = \arg\min_{\thet} \tfrac12\|\A_i\thet - \Mat b_i\|_2^2 .
\label{eq:lsq}
\end{equation}
\paragraph{Which gradient, and why zero.} The cost in \eqref{eq:lsq},
$f(\thet) := \tfrac12\|\A\thet-\Mat b\|_2^2$, is a function of the parameter
vector $\thet$ alone --- $\A$ and $\Mat b$ are fixed data. Its gradient with
respect to $\thet$ is
\begin{equation}
\nabla_{\thet} f(\thet) = \A^{\!\top}(\A\thet - \Mat b) = \A^{\!\top}\A\,\thet - \A^{\!\top}\Mat b .
\label{eq:grad}
\end{equation}
$f$ is convex in $\thet$: its Hessian (second derivative) is $\A^{\!\top}\A$,
\textbf{positive semi-definite (PSD)} --- a symmetric matrix $\Mat M$ is PSD if
and only if $\Mat x^{\!\top}\Mat M\Mat x \ge 0$ for every $\Mat x$, equivalently
if and only if every eigenvalue of $\Mat M$ is $\ge 0$ --- by \eqref{eq:psd-def}
below, so $f$ has no local minima that are
not global. For a convex, differentiable function this is exactly the
condition that identifies a minimiser: $\nabla_{\thet} f(\hat\thet)=0$ is
necessary (any point with nonzero gradient has a downhill direction, so it
cannot minimise) and, by convexity, also sufficient (a stationary point of a
convex function is a global one). Setting \eqref{eq:grad} to zero therefore
gives the \textbf{normal equations}
\src{Gauss 1809; Legendre 1805; see Golub \& Van Loan, \emph{Matrix
Computations} 4th ed., §5.3}
\begin{equation}
\underbrace{\A_i^{\!\top}\A_i}_{\textstyle \Mat G_i}\,\hat\thet_i
  = \A_i^{\!\top}\Mat b_i .
\label{eq:normal}
\end{equation}
$\Mat G_i := \A_i^{\!\top}\A_i \in \mathbb R^{n\times n}$ is the \textbf{Gram
matrix}. It is symmetric and positive semi-definite (PSD: sign fixed, can be zero): $\Mat
G$ qualifies because
\begin{equation}
\Mat x^{\!\top}\Mat G\Mat x = \|\A\Mat x\|^2 \ge 0 .
\label{eq:psd-def}
\end{equation}
(PSD, not the stronger positive \emph{definite}, because $\Mat x^{\!\top}\Mat
G\Mat x=0$ is possible for $\Mat x\ne 0$ whenever $\A$ has a kernel --- exactly
the degenerate directions of \S\ref{sec:kernel}.)

\paragraph{The Gram matrix is the Hessian (second derivative).} Differentiating \eqref{eq:lsq}
twice,
\begin{equation}
\nabla^2_{\thet}\, \tfrac12\|\A\thet-\Mat b\|^2 = \A^{\!\top}\A = \Mat G .
\label{eq:hessian}
\end{equation}
This identity is why the sloppy-model literature, which speaks of the Hessian
(second derivative) of the cost, and the numerical-linear-algebra literature, which speaks of the Gram
matrix, are describing the same object here \src{Gutenkunst et al.\ 2007, PLoS
Comput.\ Biol.; Transtrum et al.\ 2015, J.~Chem.\ Phys.}. For a nonlinear model
the equality is only approximate (Gauss--Newton); for \eqref{eq:ode} it is exact.

\section{Dynamical equivalence and non-identifiability}
\label{sec:kernel}

\subsection{Setup}

Generate the trajectory by explicit Euler from a fixed $v_i(0)$ (unit time
step, matching the stacking in \eqref{eq:ls}). Row $t$ of \eqref{eq:ls} is
the per-step constraint that maps $v_i(t)\mapsto v_i(t+1)$ given the
parameters and the already-known presynaptic activity at $t$,
\begin{equation}
\Mat a_i(t)^{\!\top}\thet_i = b_i(t), \qquad
\Mat a_i(t) := \big(v_i(t{+}1)-v_i(t),\ -1,\ -\mathrm{ReLU}(v_{j_1}(t)),\dots\big) ,
\label{eq:row}
\end{equation}
so $\Mat a_i(t)$ is literally row $t$ of $\A_i$, and $\Mat
a_i(t)^{\!\top}\boldsymbol\eta$ is the $t$-th entry of $\A_i\boldsymbol\eta$.
This is the object both definitions below are built from.

\subsection{Definitions}

\paragraph{Non-identifiability (a degenerate direction).} A vector
$\boldsymbol\eta\neq\Mat 0$ is a \emph{degenerate direction} for neuron $i$ if
$\A_i\boldsymbol\eta=\Mat 0$, i.e.\ $\boldsymbol\eta\in\ker\A_i$. If $\rank
\A_i<n$ there exist $k:=n-\rank\A_i$ such linearly independent vectors
$\boldsymbol\eta_1,\dots,\boldsymbol\eta_k$ (rank--nullity: $k=\dim\ker\A_i$),
and for every $\lambda_1,\dots,\lambda_k\in\mathbb R$,
\begin{equation}
\A_i\Big(\thet_i + \textstyle\sum_{l=1}^k \lambda_l\boldsymbol\eta_l\Big) = \A_i\thet_i = \Mat b_i ,
\label{eq:kernel}
\end{equation}
so the data cannot distinguish $\thet_i$ from $\thet_i+\sum_l\lambda_l\boldsymbol\eta_l$:
a continuum ($\lambda_l\in\mathbb R$) in each of $k$ independent directions,
an affine solution set of dimension $k$. Note
\begin{equation}
\ker \A_i = \ker \Mat G_i ,
\label{eq:kereq}
\end{equation}
since $\Mat G\Mat x = \Mat 0 \Rightarrow \Mat x^{\!\top}\Mat G\Mat x =
\|\A\Mat x\|^2 = 0 \Rightarrow \A\Mat x = \Mat 0$, so one may look for exact
degeneracy in either matrix. This is \emph{structural} (a.k.a.\ \emph{a
priori}) non-identifiability \src{Bellman \& Åström 1970, Math.\ Biosci.;
Raue et al.\ 2009, Bioinformatics}.

\paragraph{Dynamical equivalence.} Define, for each neuron $k$, the one-step
update it actually simulates by,
\begin{equation}
\mathrm{Step}_k[\thet_k](t,\Mat u) := u_k + \frac{1}{\tau_k}\Big(-u_k +
V_k^{\mathrm{rest}} + \textstyle\sum_j W_{kj}\,\mathrm{ReLU}(u_j) + I_k(t)\Big) ,
\qquad \text{defined iff } \tau_k\neq0,
\label{eq:step}
\end{equation}
so that $v(t{+}1)_k := \mathrm{Step}_k[\thet_k](t,v(t))$ for every $k$ is
\emph{how $v$ itself was generated}, recursively, from $v(0)$. Now fix a
degenerate-direction candidate $\boldsymbol\eta\neq\Mat 0$ and a scalar
$\lambda\neq0$, and define a \emph{second, independently recursed} sequence
$w$ by
\begin{equation}
w(0):=v(0), \qquad
w(t{+}1)_i := \mathrm{Step}_i[\thet_i+\lambda\boldsymbol\eta](t,w(t)), \qquad
w(t{+}1)_k := \mathrm{Step}_k[\thet_k](t,w(t))\ \ (k\neq i) ,
\label{eq:wroll}
\end{equation}
assumed well-defined, i.e.\ $\tau_i+\lambda\eta_\tau\neq0$ throughout (only
neuron $i$'s parameters are perturbed). Note that $w$'s recursion at every
step consults \emph{its own history} $w(t)$, not $v(t)$ --- it is not
notation for the same object as $v$, only initialised the same way. The
perturbation is \emph{(exactly) dynamically equivalent} if
\begin{equation}
w(t) = v(t) \qquad \text{for every sampled } t.
\label{eq:dyneq-def}
\end{equation}

\subsection{The equivalence statement}

\begin{theorem}[Non-identifiability $\iff$ dynamical equivalence, for this
ODE class]
\label{thm:equiv}
Restrict to neurons governed by the linear-in-parameters, Euler-integrated
ODE class of \eqref{eq:ode}--\eqref{eq:ls}, with per-step update
\eqref{eq:step} --- this restriction is load-bearing (see the Remark below)
and does not hold for a generic nonlinear recursion. Within that
class, fix $\thet_i$, the trajectory $v$ generated by \eqref{eq:step} from
$v(0)$, $\boldsymbol\eta\neq\Mat 0$, and $\lambda\neq 0$ such that $w$ in
\eqref{eq:wroll} is well-defined. Then
\begin{equation}
\boldsymbol\eta\in\ker\A_i
\qquad\Longleftrightarrow\qquad
w(t)=v(t) \ \text{ for every sampled } t
\qquad (w \text{ as in \eqref{eq:wroll}, built from this } \lambda) .
\label{eq:equiv-statement}
\end{equation}
\end{theorem}

\begin{remark}[this is not a tautology]
\label{rem:nottaut}
The two sides are different kinds of object. The left, $\boldsymbol\eta\in\ker\A_i$,
is \emph{static}: the recorded row equations \eqref{eq:row}, evaluated on the
\emph{already-recorded} $v$, still balance under
$\thet_i\mapsto\thet_i+\lambda\boldsymbol\eta$. The right, $\forall t,\
w(t)=v(t)$, is \emph{dynamic}: a statement about $w$, the independently
recursed sequence \eqref{eq:wroll}, which at every step consults its own
history, not $v$'s. Nothing in the definition of $\ker\A_i$ says a
freshly-recursed sequence must reproduce $v$ --- for a generic nonlinear
recursion a one-step degeneracy does \emph{not} keep the rollout locked to a
reference trajectory under a \emph{finite} parameter change; perturbations
generically compound rather than cancel. Hence $(\Rightarrow)$ needs
induction, and $(\Leftarrow)$ is a genuine reduction from a global dynamical
fact to a per-row algebraic one. The content of the theorem is that the two
coincide \emph{for this ODE class} --- a structural fact about this model
class, not a general property of recursive systems.

\emph{What must be proved to rule out the tautology, and how Lean settles it.}
``Tautology'' would mean the two sides are the same statement, so the
equivalence would hold \emph{with no hypotheses}. The way to refute that is
therefore to exhibit a hypothesis that is genuinely load-bearing --- one
whose removal makes the theorem false. The candidate is $\lambda\neq0$, and
\S\ref{sec:lean} discharges it mechanically: the companion Lean theorem
\texttt{dyneq\_trivial\_at\_zero} proves that at $\lambda=0$ the right-hand
side holds \emph{unconditionally} (the perturbation is inert, so $w\equiv
v$ whatever $\boldsymbol\eta$ is). If $(\Leftarrow)$ held without
$\lambda\neq0$, composing the two would make $\boldsymbol\eta\in\ker\A_i$
unconditionally true --- false for any $\A_i$ of full rank. So
$\lambda\neq0$ cannot be dropped, the right-hand side carries information
the left does not, and the equivalence is a theorem rather than a
restatement.
\end{remark}

\subsection{Demonstration}

\begin{proof}[Proof of $(\Rightarrow)$: a degenerate direction leaves the rollout invariant]

\emph{Base case.}
\begin{equation}
w(0) = v(0) \qquad \text{by construction.}
\label{eq:base}
\end{equation}

\emph{Inductive hypothesis.}
\begin{equation}
w(t) = v(t) \qquad \text{(the \emph{whole} population state).}
\label{eq:hyp}
\end{equation}
This must be the whole state, not just neuron $i$'s: $\Mat a_i(t)$ reads off
presynaptic $\mathrm{ReLU}(v_j(t))$ too, so \eqref{eq:hyp} is needed for every
neuron feeding $i$, not just $i$ itself.

\emph{Inductive step (show \eqref{eq:hyp} $\Rightarrow$ $w(t{+}1)=v(t{+}1)$).}
The perturbed row-$t$ equation is
\begin{equation}
\Mat a_i(t;w)^{\!\top}(\thet_i+\lambda\boldsymbol\eta) = b_i(t), \qquad
\Mat a_i(t;w) := \big(w_i(t{+}1)-w_i(t),\ -1,\ -\mathrm{ReLU}(w_{j_1}(t)),\dots\big) .
\label{eq:rowpert}
\end{equation}
Substituting \eqref{eq:hyp} collapses every entry but the first onto the true,
known value,
\begin{equation}
\Mat a_i(t;w) = \big(w_i(t{+}1)-v_i(t),\ -1,\ -\mathrm{ReLU}(v_{j_1}(t)),\dots\big) ,
\end{equation}
so the only unknown left is $w_i(t+1)$, in the first entry. Set
$\delta := w_i(t+1)-v_i(t+1)$. Subtracting the true row-$t$ equation
$\Mat a_i(t)^{\!\top}\thet_i=b_i(t)$ from \eqref{eq:rowpert} isolates
\begin{equation}
(\tau_i+\lambda\eta_\tau)\,\delta \;+\; \lambda\,\Mat a_i(t)^{\!\top}\boldsymbol\eta \;=\; 0 ,
\label{eq:induct}
\end{equation}
where $\eta_\tau$ is the $\tau_i$-component of $\boldsymbol\eta$. The second
term is $\lambda$ times \emph{exactly} the kernel condition --- $\Mat
a_i(t)^{\!\top}\boldsymbol\eta$ is the $t$-th entry of $\A_i\boldsymbol\eta$,
which is $\Mat 0$ because $\boldsymbol\eta\in\ker\A_i$ --- so it vanishes and
\eqref{eq:induct} reduces to
\begin{equation}
(\tau_i+\lambda\eta_\tau)\,\delta=0 .
\label{eq:induct-reduced}
\end{equation}
Provided the perturbed leak coefficient stays nonzero,
$\tau_i+\lambda\eta_\tau\neq 0$, \eqref{eq:induct-reduced} forces $\delta=0$,
i.e.\ $w_i(t+1)=v_i(t+1)$. Every other neuron $k\neq i$ sharing the population
state has $\thet_k$ untouched and inputs $w(t)=v(t)$ by \eqref{eq:hyp}, so
$w_k(t+1)=v_k(t+1)$ trivially, by its own unperturbed update rule. Together,
\begin{equation}
w(t+1) = v(t+1) ,
\end{equation}
which is \eqref{eq:hyp} one step later --- the inductive step is complete.

\emph{Conclusion.} By \eqref{eq:base} and the inductive step, $w(t)=v(t)$ for
every sampled $t$: a null direction leaves the rollout invariant, not merely
the one-shot fit. This proves $(\Rightarrow)$.
\end{proof}

\begin{proof}[Proof of $(\Leftarrow)$: an exactly dynamics-preserving
perturbation must be a degenerate direction]

No induction is needed here, since the hypothesis already gives the
trajectory match at every $t$ directly, rather than building it up step by
step. Suppose $w(t)=v(t)$ for \emph{every} sampled $t$, under the fixed
$\lambda\neq 0$. Unlike the inductive step above, $w(t{+}1)=v(t{+}1)$ is now
already given (it is part of ``for every sampled $t$''), so substituting into
the perturbed row-$t$ equation \eqref{eq:rowpert} collapses $\Mat a_i(t;w)$
fully to $\Mat a_i(t)$, including its first entry:
\begin{equation}
\Mat a_i(t)^{\!\top}(\thet_i+\lambda\boldsymbol\eta) = b_i(t) ,
\label{eq:conv-sub}
\end{equation}
and subtracting the true row-$t$ equation \eqref{eq:row},
$\Mat a_i(t)^{\!\top}\thet_i=b_i(t)$, cancels every $\thet_i$ and $b_i(t)$ term,
leaving
\begin{equation}
\lambda\,\Mat a_i(t)^{\!\top}\boldsymbol\eta = 0 .
\label{eq:conv-diff}
\end{equation}
Since $\lambda\neq 0$,
\begin{equation}
\Mat a_i(t)^{\!\top}\boldsymbol\eta = 0 \qquad \text{for every sampled } t,
\label{eq:conv-concl}
\end{equation}
i.e.\ $\boldsymbol\eta\in\ker\A_i$. This proves $(\Leftarrow)$.
\end{proof}

\subsection{Lean formalization}
\label{sec:lean}

Theorem~\ref{thm:equiv} has been machine-checked in Lean~4 (with Mathlib);
the development is in \texttt{lean\_proof/equiv\_proof/EquivProof/Basic.lean}.
The top-level statement is short:
\begin{verbatim}
theorem nonident_iff_dyneq
    (n : Nat) (tau Vrest : Fin n -> Real) (W : Fin n -> Fin n -> Real)
    (Iext : Fin n -> Nat -> Real) (i : Fin n) (eta_tau eta_V : Real)
    (eta_W : Fin n -> Real) (lam : Real) (v0 : State n) (T : Nat)
    (htau  : forall k, tau k <> 0)          -- Euler step well-defined
    (htau' : tau i + lam * eta_tau <> 0)    -- perturbed step well-defined
    (hlam  : lam <> 0) :                    -- LOAD-BEARING, see below
    IsDegenerate n tau Vrest W Iext i eta_tau eta_V eta_W v0 T
      <-> IsDynEquiv n tau Vrest W Iext i eta_tau eta_V eta_W lam v0 T
\end{verbatim}
(The Lean source uses Unicode identifiers --- $\mathbb N$, $\tau$,
$\boldsymbol\eta$, $\neq$, $\leftrightarrow$ --- transliterated to ASCII
here for typesetting only.)

\paragraph{The statement is short because the definitions carry the content
--- so the definitions, not the theorem line, are what must be audited.}
This is the single most important thing to understand about reading a
formalization, and it cuts against the natural reading. A machine-checked
proof guarantees only that the \emph{stated} proposition follows from the
\emph{stated} definitions; it cannot tell you that those definitions say
what you meant. Lean checks the proof; \emph{the reader} must check the
specification. Since \texttt{nonident\_iff\_dyneq} is a one-line claim about
two named predicates, essentially all of its meaning lives in those
predicates and in the recursions they mention. They are therefore given
here in full, so that the specification is auditable without opening the
source:
\begin{verbatim}
-- Population state: one real number per neuron.
abbrev State (n : Nat) := Fin n -> Real

-- One neuron k's explicit-Euler update, eq:step. Note it reads the FULL
-- population state u, not just u k: the presynaptic ReLU terms are in the sum.
def stepOne n tau Vrest W Iext (k : Fin n) (t : Nat) (u : State n) : Real :=
  u k + (1 / tau k) * (-(u k) + Vrest k + (sum j, W k j * relu (u j)) + Iext k t)

def stepAll n tau Vrest W Iext t (u : State n) : State n :=
  fun k => stepOne n tau Vrest W Iext k t u

-- v : the TRUE trajectory, defined BY the recursion, from v0.
def traj n tau Vrest W Iext (v0 : State n) : Nat -> State n
  | 0     => v0
  | t + 1 => stepAll n tau Vrest W Iext t (traj n tau Vrest W Iext v0 t)

-- Only neuron i's parameters are perturbed: `Function.update f i val` is
-- "the function f, except it returns val at index i".
def tau_pert   n tau   i eta_tau lam := Function.update tau   i (tau i + lam * eta_tau)
def Vrest_pert n Vrest i eta_V   lam := Function.update Vrest i (Vrest i + lam * eta_V)
def W_pert     n W     i eta_W   lam := Function.update W i (fun j => W i j + lam * eta_W j)

-- w : a SECOND, INDEPENDENT recursion. At step t it consults its OWN
-- history (trajPert ... t), never traj's. Only v0 is shared.
def trajPert n tau Vrest W Iext i eta_tau eta_V eta_W lam (v0 : State n) : Nat -> State n
  | 0     => v0
  | t + 1 => stepAll n (tau_pert n tau i eta_tau lam) (Vrest_pert n Vrest i eta_V lam)
               (W_pert n W i eta_W lam) Iext t
               (trajPert n tau Vrest W Iext i eta_tau eta_V eta_W lam v0 t)

-- LHS: eta in ker A_i. The row equation eq:row, dotted with eta, vanishing
-- at every sampled t, evaluated on the ALREADY-RECORDED trajectory traj.
def IsDegenerate ... (T : Nat) : Prop :=
  forall t < T,
    eta_tau * (traj ... (t+1) i - traj ... t i)
      - eta_V - sum j, eta_W j * relu (traj ... t j) = 0

-- RHS: dynamical equivalence. An equality of FUNCTIONS State n, hence a
-- whole-population claim, at every sampled t.
def IsDynEquiv ... (T : Nat) : Prop :=
  forall t <= T, trajPert ... t = traj ... t
\end{verbatim}
(Argument lists are elided as \texttt{...} where they merely repeat;
\texttt{sum j,} is Lean's \texttt{$\sum$ j,} over all neurons.)

\paragraph{How the Flyvis ODE \eqref{eq:ode} becomes the Lean text.} Nothing
above mentions an ODE, so it is worth tracing the three steps that carry
\eqref{eq:ode} into \texttt{stepOne} and then into \texttt{IsDegenerate}.
Each step is an identity, not an approximation of the modelling kind: the
only discretisation is Euler, and it is the \emph{same} Euler used to
generate the data.

\emph{Step 1: \eqref{eq:ode} $\to$ \texttt{stepOne}.} Explicit Euler on
\eqref{eq:ode} with unit step, solved for $v_i(t{+}1)$, is
\begin{equation}
v_i(t{+}1) \;=\; \underbrace{v_i(t)}_{\texttt{u k}}
  + \frac{1}{\tau_i}\Big(\underbrace{-v_i(t)}_{\texttt{-(u k)}}
  + \underbrace{V_i^{\mathrm{rest}}}_{\texttt{Vrest k}}
  + \underbrace{\textstyle\sum_j W_{ij}\mathrm{ReLU}(v_j(t))}_{\texttt{sum j, W k j * relu (u j)}}
  + \underbrace{I_i(t)}_{\texttt{Iext k t}}\Big),
\label{eq:ode-to-step}
\end{equation}
which is \texttt{stepOne} verbatim, brace for brace. Iterating it is
\texttt{traj}. So \texttt{traj} is not a fresh modelling choice: it is
\eqref{eq:ode} under the integrator, and the $1/\tau_i$ is why
\texttt{htau : forall k, tau k <> 0} appears among the hypotheses.

\emph{Step 2: \eqref{eq:ode} is linear in $\thet_i$.} Multiplying
\eqref{eq:ode-to-step} by $\tau_i$ and collecting the unknowns
$\thet_i=(\tau_i,V_i^{\mathrm{rest}},W_{ij_1},\dots)$ on the left gives
exactly the row equation \eqref{eq:row}:
\begin{equation}
\underbrace{\tau_i\big(v_i(t{+}1)-v_i(t)\big) - V_i^{\mathrm{rest}}
  - \textstyle\sum_j W_{ij}\mathrm{ReLU}(v_j(t))}_{\textstyle \Mat a_i(t)^{\!\top}\thet_i}
\;=\; \underbrace{-v_i(t)+I_i(t)}_{\textstyle b_i(t)} .
\label{eq:ode-to-row}
\end{equation}
The nonlinearity $\mathrm{ReLU}$ sits entirely in the coefficients $\Mat
a_i(t)$, never on the unknowns --- this is the ``no linearisation has
occurred'' point of \S\ref{sec:kernel}, and it is what makes a purely
algebraic kernel condition legitimate.

\emph{Step 3: \eqref{eq:ode-to-row} $\to$ \texttt{IsDegenerate}.} Replacing
$\thet_i$ by $\boldsymbol\eta=(\eta_\tau,\eta_V,\eta_W)$ in the left-hand
side of \eqref{eq:ode-to-row} --- i.e.\ forming $\Mat
a_i(t)^{\!\top}\boldsymbol\eta$, the $t$-th entry of $\A_i\boldsymbol\eta$
--- and asking it to vanish for every sampled $t$ gives, term for term,
\begin{center}
\begin{tabular}{@{}ll@{}}
\toprule
$\Mat a_i(t)^{\!\top}\boldsymbol\eta = 0$ & \texttt{IsDegenerate} \\
\midrule
$\eta_\tau\big(v_i(t{+}1)-v_i(t)\big)$
  & \texttt{eta\_tau * (traj ... (t+1) i - traj ... t i)} \\
$-\,\eta_V$
  & \texttt{- eta\_V} \\
$-\sum_j \eta_{W_j}\mathrm{ReLU}(v_j(t))$
  & \texttt{- sum j, eta\_W j * relu (traj ... t j)} \\
$=0$ for every sampled $t$
  & \texttt{= 0}, under \texttt{forall t < T} \\
\bottomrule
\end{tabular}
\end{center}
Note what is \emph{absent}: $b_i(t)$ does not appear, because it cancels
when the perturbed row equation is subtracted from the true one
(\eqref{eq:conv-diff}); and \texttt{traj} appears throughout, so the
condition is read off the \emph{recorded} trajectory, exactly as
\eqref{eq:kernel} requires. \texttt{IsDynEquiv}, by contrast, mentions
\texttt{trajPert} --- the re-simulated rollout --- which is the whole
asymmetry the theorem is about.

\subsection{The proof}
\label{sec:leanproof}

The proof is 79 lines (L86--164 of \texttt{EquivProof/Basic.lean}) and is
reproduced below, with only the repeated argument lists elided --- every
tactic invocation is shown. Line numbers are the file's own, so the gaps
mark elided continuation lines. A \emph{tactic} is an instruction that
transforms the current goal; the proof is read top to bottom as a sequence
of such transformations. Comments (\texttt{--}) are in the source.

\begin{verbatim}
 86  constructor
 87  · -- (=>): Base case / Inductive hypothesis / Inductive step / Conclusion.
 88    intro hdeg t ht
 89    induction t with
 90    | zero => rfl
 91    | succ t ih =>
 92      have ht' : t < T := by omega
 93      have ihfull : trajPert n tau Vrest W Iext i eta_tau eta_V eta_W lam v0 t
 94          = traj n tau Vrest W Iext v0 t := ih (le_of_lt ht')
 95      have hdeg_t := hdeg t ht'
 96      -- Unfold `traj` at `t+1`: this is exactly the row equation (eq:row).
 97      have hunfold : traj n tau Vrest W Iext v0 (t + 1) i
 98          = traj n tau Vrest W Iext v0 t i + (1 / tau i) *
 99              (-(traj n tau Vrest W Iext v0 t i) + Vrest i
100                + (sum j, W i j * relu (traj n tau Vrest W Iext v0 t j))
101                + Iext i t) := rfl
102      funext k
103      show stepAll n (tau_pert ...) (Vrest_pert ...) (W_pert ...) Iext t
104              (trajPert ... v0 t) k
105          = stepAll n tau Vrest W Iext t (traj n tau Vrest W Iext v0 t) k
106      rw [ihfull]
107      unfold stepAll
108      rcases eq_or_ne k i with hk | hk
109      · subst hk
110        unfold stepOne tau_pert Vrest_pert W_pert
111        simp only [Function.update_self]
112        have hsum : sum j, (W k j + lam * eta_W j) * relu (traj ... t j)
113            = (sum j, W k j * relu (traj ... t j))
114              + lam * sum j, eta_W j * relu (traj ... t j) := by
115          rw [Finset.mul_sum, <- Finset.sum_add_distrib]
116          exact Finset.sum_congr rfl (fun j _ => by ring)
117        rw [hsum, add_right_inj]
118        rw [hunfold] at hdeg_t
119        have htau'' : tau k + eta_tau * lam <> 0 := by rw [mul_comm]; exact htau'
120        field_simp [htau k, htau', htau''] at hdeg_t |-
121        linear_combination (-lam) * hdeg_t
122      · unfold stepOne
123        have h1 : tau_pert   n tau   i eta_tau lam k = tau k   := Function.update_of_ne hk _ _
124        have h2 : Vrest_pert n Vrest i eta_V   lam k = Vrest k := Function.update_of_ne hk _ _
125        have h3 : W_pert     n W     i eta_W   lam k = W k     := Function.update_of_ne hk _ _
126        rw [h1, h2, h3]
127  · -- (<=): direct substitution, no induction.
128    intro hdyn t ht
129    have h1 := hdyn (t + 1) (by omega)
130    have h0 := hdyn t (by omega)
131    have hunfold : traj n tau Vrest W Iext v0 (t + 1) i
132        = ...                                            := rfl
135    have hunfold' : trajPert ... (t + 1) i
136        = trajPert ... t i + (1 / (tau i + lam * eta_tau)) *
137            (-(trajPert ... t i) + (Vrest i + lam * eta_V)
138              + (sum j, (W i j + lam * eta_W j) * relu (trajPert ... t j))
139              + Iext i t) := by
141      show stepOne n (tau_pert ...) (Vrest_pert ...) (W_pert ...) Iext i t
142              (trajPert ... t) = _
144      unfold stepOne tau_pert Vrest_pert W_pert
145      simp only [Function.update_self]
146    rw [h0] at hunfold'
147    rw [h1] at hunfold'
148    have hsum : sum j, (W i j + lam * eta_W j) * relu (traj ... t j)
149        = (sum j, W i j * relu (traj ... t j))
150          + lam * sum j, eta_W j * relu (traj ... t j) := by
151      rw [Finset.mul_sum, <- Finset.sum_add_distrib]
152      exact Finset.sum_congr rfl (fun j _ => by ring)
153    rw [hsum] at hunfold'
154    rw [hunfold] at hunfold'
155    rw [add_right_inj] at hunfold'
156    rw [hunfold]
157    field_simp [htau i, htau'] at hunfold' |-
158    have hcancel : lam * (eta_tau * (traj ... t i * tau i
159          + (-traj ... t i + Vrest i + sum x, W i x * relu (traj ... t x) + Iext i t)
161          - traj ... t i * tau i))
162        = lam * (tau i * eta_V + tau i * sum x, eta_W x * relu (traj ... t x)) := by
163      linear_combination hunfold'
164    linear_combination mul_left_cancel<>0 hlam hcancel
\end{verbatim}
(Repeated argument lists are elided as \texttt{...}; \texttt{sum},
\texttt{<-}, \texttt{<>}, \texttt{|-} and \texttt{mul\_left\_cancel<>0} stand
for the source's $\sum$, $\leftarrow$, $\neq$, $\vdash$ and
\texttt{mul\_left\_cancel}$_0$.)

\subsection{The proof, line by line}
\label{sec:leanlines}

\paragraph{Splitting the equivalence (L86--88).} \texttt{constructor} splits
the goal $A\leftrightarrow B$ into the two implications; the \texttt{·}
bullets focus on each in turn. \texttt{intro hdeg t ht} \emph{assumes} the
antecedent --- it moves \texttt{IsDegenerate} into the context as a named
hypothesis \texttt{hdeg} and fixes an arbitrary time \texttt{t} with its
bound \texttt{ht}. It does not prove anything; in Lean a hypothesis is an
argument you are handed.

\paragraph{The induction (L89--91).} \texttt{induction t with | zero => ...
| succ t ih => ...} is the r\'ecurrence of \S\ref{sec:kernel}, with
\texttt{zero} the initialisation, \texttt{succ} the h\'er\'edit\'e, and
\texttt{ih} the hypoth\`ese de r\'ecurrence. The base case is closed by
\texttt{rfl} --- ``true by definition'': both rollouts are \emph{defined} to
be $v_0$ at $t=0$, so there is literally nothing to prove.

\paragraph{Instantiating the hypotheses (L92--95).} \texttt{omega} is a
decision procedure for linear arithmetic over $\mathbb N$; it discharges the
index bookkeeping between \texttt{t < T} and \texttt{t + 1 <= T}.
\texttt{ih (le\_of\_lt ht')} and \texttt{hdeg t ht'} are ordinary
\emph{function applications}: a universally quantified hypothesis is a
function from (a time, a proof of its bound) to the instance at that time.

\paragraph{The bridge to the row equation (L97--101).} \texttt{hunfold}
states that $v_i(t{+}1)$ equals one Euler step applied to $v(t)$, and proves
it with \texttt{rfl}. This is the load-bearing structural point: unfolding
the recursion is true \emph{by definition}, and what it unfolds to
\emph{is} the row equation \eqref{eq:row}. The link between the simulator and
the linear algebra costs nothing because it is definitional.

\paragraph{Reducing to one neuron at a time (L102--105).} \texttt{funext k}
is function extensionality: to prove two \emph{functions} equal, prove them
equal at each argument. Since the goal is an equality of population states,
\texttt{k} now ranges over \emph{all} neurons --- the whole-population
caveat of \S\ref{sec:kernel} appears here as a proof obligation, not as a
remark. \texttt{show} restates the goal in a definitionally equal form
(one rollout step on each side).

\paragraph{The case split (L106--108).} After \texttt{rw [ihfull]} replaces
$w(t)$ by $v(t)$ --- this is the inductive hypothesis being used ---
\texttt{rcases eq\_or\_ne k i with hk | hk} splits on whether the neuron
being updated is the perturbed one. This is exactly the note's split between
neuron $i$ and ``every other neuron $k\neq i$''.

\paragraph{Case $k=i$: consuming the degeneracy (L109--121).} After
\texttt{subst} and unfolding, \texttt{Function.update\_self} supplies the
perturbed parameters at the updated index. \texttt{hsum} splits the
perturbed synaptic sum into (true part) $+\ \lambda\cdot(\boldsymbol\eta$
part$)$ --- the algebraic content of ``linear in the parameters'';
note \texttt{<-\,Finset.sum\_add\_distrib} is applied \emph{backwards}, to
recombine two sums into one before closing termwise with \texttt{ring}.
\texttt{add\_right\_inj} then cancels the common leading term (despite its
name it cancels on the \emph{left}: \texttt{a + b = a + c <-> b = c}), and
\texttt{field\_simp} clears the $1/\tau$ denominators, which is why the
nonvanishing hypotheses are passed to it. The last line is the punchline:
\begin{center}
\texttt{linear\_combination (-lam) * hdeg\_t}
\end{center}
proves the goal by exhibiting $\text{goal}-(-\lambda)\cdot(\text{degeneracy
hypothesis})$ as a ring identity. Mathematically: \emph{the perturbed update
agrees with the true one precisely because $\Mat
a_i(t)^{\!\top}\boldsymbol\eta=0$, scaled by $\lambda$}. This single line is
\eqref{eq:induct}$\to$\eqref{eq:induct-reduced}.

\paragraph{Case $k\neq i$: the other neurons (L122--126).}
\texttt{Function.update\_of\_ne} says a function is unchanged away from the
updated index, so the untouched neurons have literally the same parameters
and hence the same update. Three rewrites close it. Trivial --- but the
formalization \emph{forces} it to be discharged, which is the note's point
that the induction hypothesis must cover the whole population.

\paragraph{The converse: no induction (L127--130).} \texttt{hdyn} is applied
at $t$ and at $t{+}1$ to extract the two trajectory agreements needed. There
is no \texttt{induction} tactic anywhere in this branch, and that is the
structural asymmetry: agreement at \emph{every} time is already a
hypothesis, so nothing has to be built up step by step.

\paragraph{Collapsing $w$ onto $v$ (L146--147).} \texttt{rw [h0]} and
\texttt{rw [h1]} substitute the dynamical-equivalence equalities into the
perturbed update equation, rewriting it entirely in terms of $v$. This is
the note's ``substituting collapses $\Mat a_i(t;w)$ fully to $\Mat
a_i(t)$''. Because \texttt{h0} is an equality of \emph{functions}, the
rewrite reaches inside $\sum_j\cdots\mathrm{ReLU}(w_j(t))$, every
presynaptic term; had \texttt{IsDynEquiv} been stated for neuron $i$ alone
this step would fail. The whole-population choice is load-bearing here.

\paragraph{Dividing by $\lambda$ (L157--164).} With both update equations
expressed in $v$ and denominators cleared, subtracting them leaves
\texttt{hcancel}, deliberately shaped as $\lambda\cdot A=\lambda\cdot B$ ---
which is \eqref{eq:conv-diff}. Then
\begin{center}
\texttt{mul\_left\_cancel}$_0$ \texttt{hlam hcancel}
\end{center}
(signature \texttt{a <> 0 -> a * b = a * c -> b = c}) is exactly the note's
step ``since $\lambda\neq0$'', done honestly: the nonvanishing fact is
consumed as an explicit argument rather than by dividing. \textbf{This is
the only place \texttt{hlam} is used in the entire proof} --- every other
step is valid at $\lambda=0$ --- so the machine pins the whole dependence on
$\lambda\neq0$ to one identifiable inference, precisely the one
Remark~\ref{rem:nottaut} says the tautology question turns on.

\paragraph{What the axiom report certifies.} \texttt{\#print axioms} reports
\texttt{[propext, Classical.choice, Quot.sound]} for both theorems: Lean's
three standard foundations, present in essentially every Mathlib
development. The meaningful part is what is \emph{absent}. \texttt{sorry} is
Lean's ``trust me'' placeholder, and any proof containing one --- at any
depth, in any lemma it depends on --- carries \texttt{sorryAx} in this list
indelibly. Its absence certifies that no step is left unproven: the
difference between ``formalised'' and ``machine-checked''.

\section{The SVD and the eigenvalue spectrum}

\subsection{Definition}
Every $\A\in\mathbb R^{T\times n}$ factors as
\begin{equation}
\A = \Mat U \boldsymbol\Sigma \Mat V^{\!\top},\qquad
\boldsymbol\Sigma = \mathrm{diag}(\sigma_1 \ge \dots \ge \sigma_n \ge 0),
\label{eq:svd}
\end{equation}
with $\Mat U,\Mat V$ orthogonal \src{Beltrami 1873; Jordan 1874; Eckart \& Young
1936, Psychometrika; Golub \& Kahan 1965, SIAM J.~Numer.\ Anal.}. The $\sigma_k$
are the \textbf{singular values}, the columns $\Mat v_k$ of $\Mat V$ the right
singular vectors --- directions in \emph{parameter} space.

\subsection{The relation you must keep straight}
Substituting \eqref{eq:svd} into the Gram matrix,
\begin{equation}
\boxed{\;\Mat G = \A^{\!\top}\A = \Mat V\boldsymbol\Sigma^2\Mat V^{\!\top}
\quad\Longrightarrow\quad
\lambda_k(\Mat G) = \sigma_k(\A)^2 \;}
\label{eq:key}
\end{equation}
So: the \emph{eigenvalues} of the Gram matrix are the \emph{squared singular
values} of the design matrix, sharing eigenvectors. Because $\Mat G$ is
symmetric PSD its eigenvalues and singular values coincide, which is why one can
get away with saying ``singular values of the Gram matrix'' --- but the
\emph{square} is then silently in play, and that is the source of the confusion
documented in §\ref{sec:flags}. Define the normalised spectrum
\begin{equation}
\eps_k := \frac{\lambda_k(\Mat G)}{\lambda_{\max}(\Mat G)}
        = \left(\frac{\sigma_k(\A)}{\sigma_1(\A)}\right)^{\!2}
        \in [0,1].
\label{eq:relspec}
\end{equation}
The ordered list $\{\eps_k\}$ is the \textbf{eigenvalue spectrum}. A spectrum
spanning many decades with no gap is the signature of a \emph{sloppy} model
\src{Gutenkunst et al.\ 2007; Transtrum et al.\ 2015}.

\subsection{Conditioning}
The 2-norm condition number is
\begin{equation}
\cond_2(\A) = \sigma_1/\sigma_n, \qquad\text{and crucially}\qquad
\cond_2(\Mat G) = \cond_2(\A)^2 .
\label{eq:cond}
\end{equation}
\src{Golub \& Van Loan §2.6, §5.3; Higham, \emph{Accuracy and Stability of
Numerical Algorithms} 2nd ed., ch.~20}

\paragraph{The practical consequence.} IEEE double precision has
$\eps_{\mathrm{mach}} \approx 2.2\times10^{-16}$. Forming $\Mat G$ numerically
squares the condition number, so information about $\eps_k$ below
$\sim\eps_{\mathrm{mach}}$ is destroyed: computed relative eigenvalues below
$\sim10^{-16}$ are numerical noise. Working with $\A$ directly via a
backward-stable QR or SVD does \emph{not} square it, so $\sigma_k/\sigma_1$
resolves to $\sim10^{-16}$ and therefore $\eps_k = (\sigma_k/\sigma_1)^2$ is
meaningful to $\sim10^{-32}$ \src{Golub \& Van Loan §5.2 (Householder QR
backward stability); Trefethen \& Bau, \emph{Numerical Linear Algebra}, Lectures
16, 18--19}. \textbf{This is exactly the gain of a factor $10^{16}$ in
resolvable tolerance that separates the submission's method from the rebuttal's,
and it is why a $10^{-22}$ threshold is meaningful in one and not the other.}

\section{Solving: pseudo-inverse and truncation}

When $\A$ is rank-deficient, \eqref{eq:lsq} has infinitely many minimisers. The
one of smallest norm is given by the \textbf{Moore--Penrose pseudo-inverse}
\src{Moore 1920; Penrose 1955, Proc.\ Camb.\ Phil.\ Soc.; Golub \& Van Loan
§5.5.4}
\begin{equation}
\A^{+} = \Mat V\boldsymbol\Sigma^{+}\Mat U^{\!\top},\qquad
(\boldsymbol\Sigma^{+})_{kk} =
\begin{cases} 1/\sigma_k & \sigma_k > 0\\ 0 & \sigma_k = 0\end{cases},
\qquad \hat\thet = \A^{+}\Mat b .
\label{eq:pinv}
\end{equation}
In floating point ``$\sigma_k>0$'' is replaced by a threshold, giving the
\textbf{truncated SVD} (TSVD) estimator
\begin{equation}
\hat\thet^{(\eps)} = \sum_{k:\,\eps_k > \eps}
  \frac{\Mat u_k^{\!\top}\Mat b}{\sigma_k}\,\Mat v_k ,
\label{eq:tsvd}
\end{equation}
a spectral-filter regulariser \src{Hansen, \emph{Rank-Deficient and Discrete
Ill-Posed Problems}, SIAM 1998, ch.~3--4}. Its smooth cousin is Tikhonov
regularisation, $\hat\thet = (\A^{\!\top}\A + \mu\Mat I)^{-1}\A^{\!\top}\Mat b$,
which multiplies each term by $\sigma_k^2/(\sigma_k^2+\mu)$ instead of cutting
it \src{Tikhonov 1963; Tikhonov \& Arsenin 1977}. \emph{The calcium
deconvolution in Q2 is the same estimator applied in Fourier space} --- there
the ``singular values'' are the kernel's frequency response $|K(f)|$.

\subsection{Error amplification: the master formula}
\label{sec:amp}
Perturb the right-hand side, $\Mat b \to \Mat b + \delta\Mat b$. From
\eqref{eq:tsvd},
\begin{equation}
\boxed{\;\delta\hat\thet = \sum_{k:\,\eps_k>\eps}
  \frac{\Mat u_k^{\!\top}\delta\Mat b}{\sigma_k}\,\Mat v_k
\quad\Longrightarrow\quad
\|\delta\hat\thet\| \;\lesssim\; \frac{\|\delta\Mat b\|}{\sigma_{\min}^{\text{kept}}}\;}
\label{eq:amp}
\end{equation}
\src{standard; Golub \& Van Loan §5.3.8, Trefethen \& Bau Lecture 18} Read it
carefully, because the whole of Part~III follows from it:
\begin{itemize}
\item A direction that is \emph{cut} contributes exactly zero error. Truncation
      is safe.
\item A direction that is \emph{kept} contributes error inversely proportional
      to its singular value. Keeping a near-null direction is dangerous.
\item Therefore the worst case is not $\sigma_k = 0$; it is $\sigma_k$
      \emph{just above the cut}.
\end{itemize}

\section{Two perturbation theorems: why noise helps, and how much}

\subsection{Weyl: noise lifts the spectrum}
For any $\A$ and perturbation $\Mat E$,
\begin{equation}
|\sigma_k(\A+\Mat E) - \sigma_k(\A)| \;\le\; \|\Mat E\|_2 \qquad \forall k .
\label{eq:weyl}
\end{equation}
\src{Weyl 1912; Mirsky 1960; Golub \& Van Loan Cor.~8.6.2; Stewart \& Sun,
\emph{Matrix Perturbation Theory} 1990, ch.~IV} An exactly null direction
($\sigma_k = 0$) can therefore be lifted, but by \emph{at most} the size of the
perturbation. Process noise of scale $\sigma_{\text{proc}}$ enters $\A_i$
through the recorded $v_j$, so
\begin{equation}
0 \;\longrightarrow\; \sigma_k \;\lesssim\; O(\sigma_{\text{proc}}) .
\label{eq:lift}
\end{equation}
\textbf{This is the theorem behind ``per-neuron noise breaks the degenerate
directions of $\A_i$''} (\texttt{neurips.tex} L354--355). It is a real
mechanism, and \eqref{eq:weyl} is the citation the claim currently lacks.

\subsection{Marchenko--Pastur: what the lift looks like}
For $\Mat E$ with i.i.d.\ entries of variance $s^2$ and aspect ratio
$n/T \to c$, the singular values of $\Mat E/\sqrt T$ fill
$[s(1-\sqrt c),\, s(1+\sqrt c)]$ \src{Marchenko \& Pastur 1967, Mat.\ Sb.; Bai
\& Silverstein, \emph{Spectral Analysis of Large Dimensional Random Matrices}}.
So noise does not merely lift one direction; it installs a \emph{floor} of width
$O(s)$ across the spectrum. Distinguishing ``noise selectively destroyed the
columnar degeneracy'' from ``noise put a floor under everything'' is thus a real
question, not a rhetorical one --- and it is precisely the question
\texttt{reply\_all.tex} L233--237 tries to settle by column equilibration.

\subsection{Column equilibration, and what it does and does not fix}
Rescaling columns, $\A \to \A\Mat D$ with $\Mat D$ diagonal, changes the
spectrum; choosing $\Mat D$ so each column has unit norm is
\emph{equilibration} and is close to optimal for conditioning among diagonal
scalings \src{van der Sluis 1969, Numer.\ Math.}. After equilibration
\begin{equation}
\sum_k \lambda_k(\Mat G) = \tr(\Mat G) = \sum_{k} \|\A_{:,k}\|^2 = n ,
\label{eq:trace}
\end{equation}
a fixed budget. So noise \emph{cannot} raise all $\eps_k$ together: raising the
small ones must lower the large ones. That is a genuine argument and
\eqref{eq:trace} is its one-line proof. But note what it does \emph{not} give
you: it constrains the sum, not the shape, so it rules out a uniform inflation
while remaining compatible with a Marchenko--Pastur floor. The decisive evidence
is the columnar decomposition (§\ref{sec:usage}, item~\ref{it:col}), not the
trace identity.

\subsection{Correlated columns force a small singular value (closed-form, not
just observed)}
\label{sec:corr-cols}

Kristoffersen, Nielsen \& Solem (2024/25) prove a closed-form
condition-number bound for their own (Heaviside firing-interval) inverse
problem, driven by \emph{how close two firing events are in time}
\src{Kristoffersen et al., arXiv:2409.07241, Thm 3.5}: two rows of their
design matrix built from firing events separated by a small gap $h$ are
forced close together, giving $\sigma_1\ge c_1$, $\sigma_r\le c_2h/\|\cdot\|$,
hence $\cond\gtrsim h^{-1}$. Our setting has a natural twin, driven instead
by \emph{how correlated two presynaptic columns are} --- the mechanism
behind the ``time-shifted copies'' argument the Area Chair objected to in
\S\ref{sec:flags}(a): their theorem needs two \emph{rows} to be
near-duplicates (close firing times); ours needs two \emph{columns} to be
near-duplicates (correlated presynaptic neurons).

\paragraph{Setup.} Let $a_p, a_q$ be two columns of $\A_i$ corresponding to
presynaptic weights $W_{ip}, W_{iq}$ (e.g.\ two same-type, time-shifted
presynaptic neurons). Define
\begin{equation}
\alpha := \frac{\langle a_p,a_q\rangle}{\|a_p\|^2}, \qquad
\rho := \frac{\langle a_p,a_q\rangle}{\|a_p\|\,\|a_q\|}, \qquad
\eps := \|a_q-\alpha a_p\| = \|a_q\|\sqrt{1-\rho^2} .
\label{eq:corr-def}
\end{equation}
$\rho$ is the (cosine/Pearson) correlation between the two columns; $\eps$
is the part of $a_q$ \emph{not} explained by $a_p$ --- exactly ``correlated,
not identical.''

\begin{theorem}
Let $\sigma_1\ge\dots\ge\sigma_n\ge0$ be the singular values of
$\A_i\in\mathbb R^{T\times n}$. Then
\begin{equation}
\sigma_1 \ge \sqrt T, \qquad
\sigma_n \le \frac{\eps}{\sqrt{1+\alpha^2}} = \frac{\|a_q\|\sqrt{1-\rho^2}}{\sqrt{1+\alpha^2}} ,
\label{eq:corr-thm}
\end{equation}
hence
\begin{equation}
\cond(\A_i) \ge \frac{\sqrt T\,\sqrt{1+\alpha^2}}{\|a_q\|\sqrt{1-\rho^2}}
\ \xrightarrow[\ \rho\to\pm1\ ]{}\ \infty .
\label{eq:corr-cond}
\end{equation}
\end{theorem}

\begin{proof}
\emph{The residual (elementary).} By definition of $\alpha$ as the
least-squares projection coefficient, $r:=a_q-\alpha a_p \perp a_p$, so
$\|a_q\|^2=\alpha^2\|a_p\|^2+\|r\|^2$ (Pythagoras); substituting
$\alpha=\rho\|a_q\|/\|a_p\|$ gives $\|r\|=\eps$. This is the analog of
their mean-value-theorem step: it turns ``correlated but not identical''
into a residual norm, instead of a temporal gap.

\emph{Upper bound on $\sigma_n$.} Since $\A_i^{\!\top}\A_i\in\mathbb
R^{n\times n}$ has no trivial padding zeros (unlike their $A^{(i)}[A^{(i)}]^\top$,
which needed a restriction to $\mathcal N([A^{(i)}]^\top)^\perp$ to exclude
them), $\sigma_n^2=\min_{x\ne0}\langle x,\A_i^{\!\top}\A_ix\rangle/\|x\|^2$
directly, with no such restriction needed. Take
$x=\eta:=e_q-\alpha e_p$: $\A_i\eta=a_q-\alpha a_p=r$, so
$\sigma_n^2\le\eps^2/(1+\alpha^2)$.

\emph{Lower bound on $\sigma_1$.} Take $x=e_V$, the coordinate for
$V_i^{\mathrm{rest}}$. That column of $\A_i$ is exactly the constant vector
$(-1,\dots,-1)\in\mathbb R^T$ (eq.~\eqref{eq:row}), so $\|\A_ie_V\|^2=T$
exactly --- no a-priori ODE bound on an unknown state needed, unlike their
$c_1$. Hence $\sigma_1^2\ge T$.

Divide.
\end{proof}

\begin{remark}[caveat: one pair here, generalizes to $k$ correlated columns]
The theorem above handles one pair $(p,q)$. Our note's actual claim
(\S\ref{sec:flags}(a): a $(k_{i\alpha}-1)$-dimensional sloppy block from
$k_{i\alpha}$ same-type neurons) needs $k-1$ directions at once. The
\emph{construction} is immediate: if $a_{q_1},\dots,a_{q_{k-1}}$ are all
well-approximated by a common reference column $a_p$, i.e.\
$a_{q_l}=\alpha_la_p+r_l$ with $\|r_l\|=\eps_l$, then
$\eta_l:=e_{q_l}-\alpha_le_p$, $l=1,\dots,k-1$, are linearly independent
(each has a unit coefficient on its own coordinate $q_l$, all distinct),
spanning a genuine $(k-1)$-dimensional subspace $S$ --- the ``rank
$\approx1$ instead of $k$'' block, made explicit rather than asserted. The
\emph{proof} needs one step more than the pairwise case: a single test
vector bounds only the single smallest singular value, but bounding all
$k-1$ of the bottom singular values at once needs the Courant--Fischer
min-max theorem,
\begin{equation}
\sigma_{n-k+2}^2 = \min_{\dim(S')=k-1}\ \max_{x\in S',\,x\ne0}\frac{\|\A_ix\|^2}{\|x\|^2}
\ \le\ \max_{x\in S,\,x\ne0}\frac{\|\A_ix\|^2}{\|x\|^2} ,
\end{equation}
choosing $S$ itself as the particular $(k{-}1)$-dimensional subspace
plugged into the min. For $x=\sum_lc_l\eta_l\in S$, $\A_ix=\sum_lc_lr_l$,
and relating $\sum_l|c_l|$ back to $\|x\|$ only needs the $\eta_l$'s own
Gram matrix $\Mat I_{k-1}+\boldsymbol\alpha\boldsymbol\alpha^{\!\top}$ (a
rank-one perturbation of the identity, condition number controlled by
$\|\boldsymbol\alpha\|$). The upshot: all $k-1$ bottom singular values are
$O(\max_l\eps_l)$, not just $\sigma_n$ --- the construction transfers
immediately; the bound needs the subspace version of the same tool.
\end{remark}

\section{One more trap: noise in the design matrix}
\label{sec:eiv}
All of the above perturbs $\Mat b$. If the \emph{regressors} are noisy the
estimator is not merely noisy but \emph{biased}: for a scalar regressor with
signal variance $s_x^2$ and noise variance $s_\nu^2$,
\begin{equation}
\hat\beta \;\xrightarrow{\;T\to\infty\;}\; \beta\cdot\frac{s_x^2}{s_x^2+s_\nu^2}
\;<\; \beta ,
\label{eq:attenuation}
\end{equation}
the \textbf{attenuation} (errors-in-variables) bias \src{Fuller,
\emph{Measurement Error Models}, Wiley 1987, §1.1; Wooldridge, \emph{Econometric
Analysis of Cross Section and Panel Data}, §4.4.2}. This matters because the
first column of $\A_i$ in \eqref{eq:ls} is $-\dot v_i$, which is built from the
same noisy voltages as $\Mat b_i$: the regressor is endogenous, and recovery of
$\tau_i$ is attenuated by an amount that \emph{grows with process noise}. The
rebuttal script's ``increment'' reparameterisation exists to remove exactly
this, by putting only time-$t$ quantities on the right-hand side:
\begin{equation}
\Delta v_i(t) = \alpha_i\big(-v_i + I_i\big) + \beta_i
  + \textstyle\sum_j \gamma_{ij}\mathrm{ReLU}(v_j) + \text{noise},
\qquad
\tau_i = \frac{\Delta t}{\alpha_i},\;\;
W_{ij} = \frac{\gamma_{ij}}{\alpha_i}.
\label{eq:increment}
\end{equation}
Two consequences to hold on to: the division by $\alpha_i$ means \emph{all} of
neuron $i$'s weights inherit any error in $\alpha_i$; and \eqref{eq:increment}
is a \emph{different design matrix} from \eqref{eq:ls}, so its conditioning is a
different question.

\part{How these tools are used in our two documents}
\label{sec:usage}

Every line below is quoted from the working copies. I group them by what they
assert.

\section{\texttt{neurips.tex}: the claims}

\begin{enumerate}
\item \textbf{L113} ``\emph{possible when both connectivity and activity are
sufficiently rich and identifiable}''; \textbf{L212} ``\emph{essential for
parameter identifiability}''. Informal use of identifiability, in the sense of
§\ref{sec:amp}: enough directions kept, none of them near-null.

\item \textbf{L237} ``\emph{Many synaptic weight configurations give rise to
identical neural activity, making the recovery of $\hat{\Mat W}$ \dots\
fundamentally non-unique}''. This is \eqref{eq:kernel} verbatim: an affine
solution set.

\item \textbf{L277} ``\emph{even with full knowledge of the underlying
equations, $\approx 10\%$ of the circuit parameters cannot be uniquely
determined}''. The $10\%$ is $6.35 + 4.01$, i.e.\ null plus sloppy at the two
thresholds of L831--832; see the flag in §\ref{sec:flags}(c).

\item \textbf{L822} ``\emph{if the column space is not full rank the system
admits a continuous degeneracy defined by the kernel $\ker(\A_i)$}''. Exactly
\eqref{eq:kernel}--\eqref{eq:kereq}. ``Column space'' should read ``column
rank''; the kernel is a property of the domain, not the column space.

\item \textbf{L828--832}, the method paragraph, and the one that carries all the
weight: ``\emph{We solve the linear system using least squares, using float64
arithmetic, column normalization of $\A_i$. We compute the singular value
decomposition of the $(2+d_i)\times(2+d_i)$ Gram matrix $\A_i^{\!\top}\A_i$ and
use the ratio $\lambda/\lambda_{\max}$ of singular values to define null (ratio
$< 64000\times10^{-22}$) and sloppy (ratio $< 64000\times10^{-12}$) directions
\dots\ \texttt{\textbackslash citep\{Gutenkunst2007-so\}}}''. Four things at once: least squares
\eqref{eq:lsq}, equilibration \eqref{eq:trace}, the Gram spectrum
\eqref{eq:key}, and the sloppiness convention. Three problems with it are listed
in §\ref{sec:flags}.

\item \textbf{L841--859}, Fig.~12 caption: ``\emph{Per-neuron least-squares
recovery}'', ``\emph{Black points are well-identified; orange (``sloppy'') and
red \dots\ per-neuron Gram matrix and are flagged as degenerate}'',
``\emph{parameter dimensions that have non-zero inner product with a null or
sloppy}''. The colouring is a per-parameter reduction of a per-\emph{direction}
quantity: parameter $\alpha$ is flagged if $|(\Mat v_k)_\alpha| >
\texttt{null\_comp\_tol}=10^{-3}$ for some flagged $\Mat v_k$. Worth stating,
since that $10^{-3}$ is a third free tolerance nowhere named in the paper.

\item \textbf{L870} ``\emph{no degeneracies purely along the $V^{\mathrm{rest}}$
direction}'': the constant column of $\A_i$ is never null, so
$\Mat e_2 \notin \ker\A_i$.

\item \textbf{L944--958} the reduced system $\Mat H_i\Mat w_i = \Mat b_i$ with
$\ker(\Mat H_i)$, for $\tau,V^{\mathrm{rest}}$ known. A submatrix of $\A_i$;
$\ker\Mat H_i \subseteq \ker\A_i$ restricted to weight coordinates.

\item \label{it:col} \textbf{L1005, L1017--1026} the structural argument:
same-type neurons across columns give ``\emph{time-shifted copies}'', so
$k_{i\alpha}$ columns of $\Mat H_i$ have ``\emph{rank $\approx 1$ instead of
$k_{i\alpha}$}'', every sum-zero perturbation is free, and
$\dim\ker(\Mat H) = \sum_i\sum_\alpha (k_{i\alpha}-1) = 308{,}160$ ($71\%$).
\textbf{Note the ``$\approx$''}: rank-$\approx$-1 gives small-but-nonzero
$\sigma_k$, hence a \emph{sloppy} subspace, not a kernel. The AC caught precisely
this; see §\ref{sec:flags}(a).

\item \textbf{L1029--1030, L1254} least squares used as a tool rather than an
object of study (group calibration; affine fit for the $\tau$ readout).

\item \textbf{L129, L161, L354--355, L464} the headline claim in four
formulations, the sharpest being L354--355: ``\emph{Per-neuron noise breaks the
degenerate directions of $\A_i$}''. Theorem: \eqref{eq:weyl}. Caveat:
\eqref{eq:lift} bounds the lift by the noise scale, and
Marchenko--Pastur says a floor appears everywhere.

\item \textbf{L461} ``\emph{the $\ell_1$ prior on $\hat{\Mat W}$ is the only
load-bearing regularizer, consistent with the analytical degeneracy}''. The
consistency is real: $\ell_1$ selects a sparse point inside the affine set
\eqref{eq:kernel}, which is what a degenerate problem needs.

\item \textbf{L2004--2020} truncated SVD of the $T\times N$ \emph{activity}
matrix. A different matrix and a different purpose (Eckart--Young low-rank
reconstruction); do not confuse its rank-$K$ with the parameter spectrum.
\end{enumerate}

\section{\texttt{reply\_all.tex}: the responses}

\begin{enumerate}
\item \textbf{L54--56}, the AC: ``\emph{Reconcile the approximately 6.4\% null
plus 4.0\% sloppy edge dimensions with the claimed 71\% exact null space}'',
``\emph{Highly correlated or time-shifted signals do not generally constitute
identical columns and therefore do not establish an exact $(k-1)$-dimensional
kernel}'', ``\emph{appears to conflate exact non-identifiability, numerical
sloppiness, and approximate dynamical equivalence}''. The AC is right on the
mathematics: correlated $\neq$ identical, so $\sigma_k$ is small but nonzero.

\item \textbf{L59--63} our concession, and the right one: Eq.~(11) is
``\emph{an approximate or sloppy subspace and not an exact null space}''; we will
show ``\emph{the spectrum of the normalized singular values (defined as
$\lambda/\lambda_{\max}$) of the Gram matrix}'' with thresholds
$10^{-22}$/$10^{-12}$; and ``\emph{approximately 70\% of the normalized singular
values are less than \texttt{1e-6}}''. That last number is the reconciliation:
$6.4\%$ and $71\%$ are one monotone curve $\Phi(\eps)$ read at two tolerances.

\item \textbf{L90--92} the three definitions. L90: ``\emph{when the matrix $A_i$
has a non-empty kernel, or equivalently, when the singular values of the Gram
matrix $A^\top A$ are zero}'' --- correct by \eqref{eq:kereq}, phrased loosely
(see §\ref{sec:flags}(b)). L91: ``\emph{the curvature of the Hessian (equal to
the Gram matrix)}'' --- this is \eqref{eq:hessian}, and it is exactly right.
L92: dynamical equivalence, as in §2.

\item \textbf{L96--97} ``\emph{non-identifiability and dynamical equivalence are
equivalent notions}'' for Euler-integrated trajectories. Proved in §2. This is
the paragraph most worth finishing.

\item \textbf{L192--244} (and the duplicate at L505--532) the new analysis:
``\emph{take the singular values of $A_i$ by SVD and rank each weight direction
by $\eps = (\sigma_k/\sigma_1)^2$}'', ``\emph{No weights are fitted and no solver
runs, so this measures the conditioning of the inverse problem}'', and the
$6.35/4.01 \to 2.44/0.32 \to 0.00/0.00$ table. Two upgrades over the submission,
both real: SVD of $\A_i$ instead of the Gram (a $10^{16}$ gain in resolvable
$\eps$, §3.3), and an estimator-free statement, which sidesteps the endogeneity
of §\ref{sec:eiv} rather than arguing about it.

\item \textbf{L233--237} ``\emph{ruled out by column equilibration, which fixes
the sum of squared singular values, so noise can only redistribute}''. This is
\eqref{eq:trace}. Sound as far as it goes; see §5.3 for what it does not settle.

\item \textbf{L682} ``\emph{process noise \dots\ is white, so the identifying
signal sits at high frequency, exactly the band the indicator attenuates}''.
The Fourier analogue of \eqref{eq:amp}: the calcium kernel is a spectral filter
whose small ``singular values'' are the high frequencies, and deconvolution
divides by them.

\item \textbf{L835, L567} the CX per-source gain degeneracy
($W_{\cdot j}\to cW_{\cdot j}$, $e^{g_j}\to e^{g_j}/c$) and the $\mu_2$ anchor.
A \emph{multiplicative} exact symmetry --- outside the linear theory above,
but the same phenomenon: a direction the data does not see.

\item \textbf{L976--999} the definitions appendix, ending
``\emph{\%\%\% CONTINUE THIS AND SHOW THAT DYNAMICAL EQUIVALENCE = DEGENERACY}''.
§2 above is that proof.
\end{enumerate}

\section{Four places the wording does not match the computation}
\label{sec:flags}

\begin{description}
\item[(a) ``exact'' vs ``approximate'' null space.] Eq.~(11)'s $308{,}160$
directions come from ``rank $\approx 1$'' blocks. Time-shifted signals are
strongly correlated, not equal, so those directions have small nonzero
$\sigma_k$: a sloppy subspace. Already conceded at L59; the fix is to stop
calling it a kernel anywhere, including L953--958.

\item[(b) eigenvalues vs singular values.] ``Singular values of the Gram
matrix'' (L831, L90) is defensible only because $\Mat G$ is symmetric PSD, and it
hides the square in \eqref{eq:key}. State once: $\eps_k =
\lambda_k(\Mat G)/\lambda_{\max}(\Mat G) = (\sigma_k(\A)/\sigma_1(\A))^2$, and
then say which one every threshold refers to.

\item[(c) the thresholds are not the same in the two documents.] The submission
uses $64000\times10^{-22}$ and $64000\times10^{-12}$ (L831--832) --- with an
explicit factor of $T$. The rebuttal analysis uses $10^{-22}$ and $10^{-12}$
with no $T$, on the grounds that equilibration plus a \emph{relative} eigenvalue
makes the quantity $T$-invariant. Both may be right for their own pipeline, but
they are numerically different tolerances, so the $6.35/4.01$ of the new table
is not the same statistic as the paper's $6.4\%/4.0\%$ even though the numbers
coincide. Reconcile explicitly or a careful referee will ask.

\item[(d) the Gram is still formed in the recovery path.] The rebuttal's own
argument is that one must not square the condition number --- and the
\emph{spectrum} path duly uses QR${}+{}$SVD. But the \emph{recovery} path
(\texttt{fig\_lstsq\_identifiability\_noise.py} L256) computes
\texttt{eigh(Xs.T @ Xs)} and truncates at $\eps>10^{-12}$, only four decades
above the $\sim10^{-16}$ floor that squaring imposes, so retained
near-threshold eigenvalues are partly numerical. Correct in principle, and the
right way round in the spectrum path --- but §\ref{sec:cause} measures it and
finds \emph{no} effect on the recovery numbers: re-solving with QR${}+{}$SVD
reproduces the Gram result to every printed digit at all three noise levels and
five tolerances. Worth fixing for hygiene, not a live problem.
\end{description}

\part{The eigenvalue analysis and the least-squares numbers do not conflict}

\section{The measurements}

Least-squares recovery of $\Mat W$ via \eqref{eq:increment}, identity-line $R^2$,
mean\,$\pm$\,SD over five folds, alongside the degenerate-edge fraction from the
spectrum:

\begin{center}
{\footnotesize
\begin{tabular}{lccccc}
\toprule
$\sigma$ & $\Rw$ all & $\Rw$ identifiable & $\Rw$ excl.\ blow-ups
 & median $|\delta W|$ & \% degenerate \\
\midrule
$0$    & $0.833\pm0.004$ & $1.0000\pm0.0001$ & $0.833$ & $0.0000$ & $10.36$ \\
$0.05$ & $-91.7\pm168.6$ & $-92.6\pm170.3$   & $0.752\pm0.033$ & $0.0020$ & $2.76$ \\
$0.5$  & $0.935\pm0.003$ & $0.935\pm0.003$   & $0.935$ & $0.0031$ & $0.00$ \\
\bottomrule
\end{tabular}}
\end{center}

\section{Why there is no contradiction}

\textbf{The two analyses measure different things, and this is the sentence to
remember: the eigenvalue analysis counts \emph{how many} directions are
determined; the least-squares recovery reports \emph{how precisely} they are
estimated.} Formally, the spectrum gives the set
$\{k : \eps_k > \eps\}$ --- a \emph{count}. Equation~\eqref{eq:amp} gives the
error on that set --- a \emph{magnitude}, $\|\delta\Mat b\|/\sigma_k$ per
direction. Nothing forces them to move together, and here they move oppositely:
\begin{itemize}
\item \emph{Count improves with noise.} $10.36\% \to 2.76\% \to 0.00\%$
      degenerate, by Weyl \eqref{eq:weyl}. Fewer directions are lost.
\item \emph{Precision degrades with noise.} Median $|\delta W|$ goes
      $0.0000 \to 0.0020 \to 0.0031$, because $\|\delta\Mat b\|$ in
      \eqref{eq:amp} grows. This is the monotone worsening my colleague
      reported, and it is exactly what the theory predicts.
\end{itemize}
Both are true simultaneously. ``Noise breaks degeneracies'' and ``noise degrades
estimates'' are statements about the two different factors in \eqref{eq:amp} ---
the denominator and the numerator. The submission's claim concerns the
denominator; the recovery scatter shows the numerator.

\paragraph{And the aggregate reflects both.} $\Rw$ over \emph{all} edges mixes
them, which is why it is non-monotone: $0.833$ at $\sigma=0$ (capped by the
$10.36\%$ of edges that are lost outright --- note $\Rw = 1.0000$ exactly on the
rest), and $0.935$ at $\sigma=0.5$ (nothing lost, everything a little noisy).
The best recovery is at the \emph{highest} noise, which is the eigenvalue
analysis's prediction, and the noise-free deficit is entirely the null space.

\section{What is genuinely wrong at $\sigma=0.05$}
\label{sec:cause}

The $-91.7$ is not a third phenomenon; it is \eqref{eq:amp} with a badly chosen
cut. At $\sigma = 0.05$, $114$ edges out of $434{,}112$ ($0.026\%$) have
$|\hat W - W|$ up to $1272$ against true weights of median $0.012$; excluding
them gives $\Rw = 0.752$. They are isolated single edges, not clustered; $\tau$
and $V^{\mathrm{rest}}$ for the same neurons are fine, so it is not a shared
$1/\alpha_i$ factor from \eqref{eq:increment}; and the same edges are clean at
$\sigma = 0$ (max error $1.08$) and at $\sigma = 0.5$ (max error $0.18$).

\subsection{Two tempting explanations, both measured and both wrong}

The obvious reading is the third bullet of §\ref{sec:amp}: at $\sigma=0$ the
columnar directions sit far below the cut and are truncated (zero variance, some
bias); at $\sigma=0.5$ the lift \eqref{eq:lift} carries them far above it; and at
$\sigma=0.05$ they are lifted to \emph{just above} the cut, retained, and divided
by. The second candidate is flag~(d): the recovery path forms the Gram, so
$\eps_k$ near $10^{-12}$ has only four decades of headroom above the
squaring-induced floor.

Both predict the same experiment. Sweep the cut $\eps_{\mathrm{cut}}$ over
several decades, and solve twice --- once with
\texttt{eigh}$(\Mat X^{\!\top}\Mat X)$, once with QR${}+{}$SVD of $\Mat X$, which
never squares the condition number. If explanation one holds, the blow-ups track
the cut. If explanation two holds, they shrink under the unsquared solver.
Running it (\texttt{neurips\_review/lstsq\_tol\_sweep.py}, fold cv04):

\begin{center}
{\footnotesize
\begin{tabular}{lccccc}
\toprule
$\eps_{\mathrm{cut}}$ & \multicolumn{2}{c}{$\sigma=0$}
 & \multicolumn{2}{c}{$\sigma=0.05$} & $\sigma=0.5$ \\
\cmidrule(lr){2-3}\cmidrule(lr){4-5}\cmidrule(lr){6-6}
 & $\Rw$ & \# blown & $\Rw$ & \# blown & $\Rw$ (\# blown) \\
\midrule
$10^{-14}$ & $0.8274$ & $0$  & $-428.83$ & $114$ & $0.9360$ ($0$) \\
$10^{-12}$ & $0.8266$ & $0$  & $-428.83$ & $114$ & $0.9360$ ($0$) \\
$10^{-10}$ & $0.8228$ & $0$  & $-428.79$ & $112$ & $0.9360$ ($0$) \\
$10^{-8}$  & $0.7959$ & $0$  & $-428.75$ & $107$ & $0.9360$ ($0$) \\
$10^{-6}$  & $-0.0338$& $56$ & $-428.86$ & $107$ & $0.9360$ ($0$) \\
\bottomrule
\end{tabular}}

{\footnotesize \emph{The Gram and QR${+}$SVD solvers agree to every printed digit
in all $30$ cells}, so they are not shown separately.}
\end{center}

\noindent Both explanations fail:
\begin{itemize}
\item \textbf{The cut is not the cause.} At $\sigma=0.05$ the blow-ups survive
      eight decades of $\eps_{\mathrm{cut}}$, $114 \to 107$. Directions sitting
      just above the cut would have been removed by moving it.
\item \textbf{The Gram is not the cause.} \texttt{eigh}$(\Mat X^{\!\top}\Mat X)$
      and QR${+}$SVD agree exactly, everywhere. Flag~(d) remains correct as a
      general principle and as the reason the \emph{spectrum} must not be
      computed from the Gram, but it has no measurable effect on the recovery
      here.
\end{itemize}

\subsection{What it actually is}

The blow-ups survive $\eps_{\mathrm{cut}} = 10^{-6}$, so the offending directions
have $\eps_k > 10^{-6}$: they are \emph{legitimately retained} by any defensible
cut, and their small singular values are real structure, not numerical noise.
Independently, the bad edges' onset statistic has median $\eps_k \approx
6.9\times10^{-5}$, in the same range. So \eqref{eq:amp} applies with
\begin{equation}
\frac{1}{\sigma_k} = \frac{1}{\sigma_1\sqrt{\eps_k}}
\;\sim\; \frac{10^{2}}{\sigma_1} \quad\text{at}\quad \eps_k \sim 10^{-5},
\label{eq:ampsize}
\end{equation}
compounded by the $1/\alpha_i$ of \eqref{eq:increment}. That is the whole story,
and it is the master formula with nothing pathological in it:
\begin{itemize}
\item $\sigma = 0$: the numerator $\Mat u_k^{\!\top}\delta\Mat b$ is zero, so the
      amplification multiplies nothing. Error $\approx 0$ (max $1.22$). The
      $\Rw = 0.833$ deficit is pure truncation \emph{bias}.
\item $\sigma = 0.05$: numerator nonzero, $\eps_k$ still $\sim10^{-5}$, so
      $\sim10^2$ amplification times the $1/\alpha_i$ factor reaches errors of
      $10^2$--$10^3$ on the worst edges. Observed max: $1272$.
\item $\sigma = 0.5$: the numerator is $10\times$ larger, but Weyl
      \eqref{eq:weyl} has lifted those directions by more --- the spectrum's
      median onset moves $4.5\times10^{-5} \to 7.9\times10^{-4}$, cutting the
      amplification $\sim4\times$ --- and the ratio improves. Max error $2.57$,
      zero blow-ups, best $\Rw$.
\end{itemize}
Between $\sigma = 0.05$ and $\sigma = 0.5$ the denominator wins. That is the
quantitative content of ``noise helps'', and it is not monotone below
$\sigma = 0.5$: the numerator grows linearly in $\sigma$ while the lift of any
particular $\eps_k$ saturates once it is well clear of the noise floor, so there
is an intermediate regime where amplification is worst. $\sigma = 0.05$ is in it.

\paragraph{One loose end, honestly flagged.} At $\sigma = 0$ with
$\eps_{\mathrm{cut}} = 10^{-6}$, \emph{over}-truncation creates $56$ blow-ups and
sends $\Rw$ to $-0.03$ --- discarding directions makes things worse, which pure
truncation bias cannot do. The likely route is \eqref{eq:increment}: cutting
directions perturbs $\alpha_i$, and every weight of that neuron carries
$1/\alpha_i$. Not measured; it does not affect the three
$\eps_{\mathrm{cut}}=10^{-12}$ numbers we report, but it is the reason to prefer
a parameterisation that does not divide by a fitted coefficient.

\section{What to report}

\begin{enumerate}
\item Report the two quantities separately and name them: the degenerate-edge
fraction (count) and the median or inlier error (precision). Reporting only
unfiltered $\Rw$ conflates them, and at $\sigma = 0.05$ lets $114$ edges out of
$434{,}112$ set the headline.
\item Say that the noise-free deficit is bias from truncation, not
misestimation: $\Rw = 1.0000$ on identifiable edges at $\sigma = 0$ is the
cleanest single number in the whole analysis and it is not currently quoted.
\item Keep the conditioning claim to the denominator of \eqref{eq:amp}, and cite
Weyl \eqref{eq:weyl} for the lift and \eqref{eq:trace} for why it is not a
uniform inflation.
\item Do not compare per-edge $\eps_k$ from $\A_i$ against recovery from
$\Mat X_i$: different matrices (§\ref{sec:eiv}). If the two are to be shown side
by side, compute the spectrum of the matrix that was actually inverted.
\end{enumerate}

\section*{Reading list, shortest path}
Trefethen \& Bau, \emph{Numerical Linear Algebra}, Lectures 4--5 (SVD),
11 (least squares), 18--19 (conditioning, stability) --- the fastest route to
\eqref{eq:key}, \eqref{eq:cond}, \eqref{eq:amp}.
Hansen, \emph{Rank-Deficient and Discrete Ill-Posed Problems}, ch.~3--4 ---
TSVD, Tikhonov, and the bias--variance picture of Part~III.
Gutenkunst et al.\ 2007 and Transtrum et al.\ 2015 --- sloppiness, and why a
multi-decade spectrum is the norm in biology rather than a pathology.
Raue et al.\ 2009 --- structural vs practical identifiability, the distinction
the AC is asking us to make.
Stewart \& Sun, \emph{Matrix Perturbation Theory}, ch.~IV --- Weyl and
friends, for §5.

\end{document}
